\clearpage\thispagestyle{otFront}\phantomsection\label{parcours}\pdfbookmark[0]{Construire son parcours}{bm-parcours}\otPageTitre{Construire son parcours}{Une année avec des retours réguliers}\fontsize{11}{15}\selectfont

\begin{tabularx}{\linewidth}{@{}lY@{}}\toprule
Étape & Travail conseillé\\\midrule
Calculer & Chapitres 1 à 4 : relatifs, nombres rationnels, puissances et racines carrées.\\[3mm]
Modéliser & Chapitres 5 à 9 : calcul littéral, équations, ratios, pourcentages, grandeurs quotients et programmes de calcul.\\[3mm]
Interpréter & Chapitres 10 et 11 : moyenne, médiane, étendue, probabilités et fréquences.\\[3mm]
Démontrer & Chapitres 12 à 17 : droites remarquables et logique, translations, droite des milieux, Pythagore, cercle circonscrit, pyramides et cônes.\\[3mm]
Programmer & Chapitre 18, à reprendre pendant l’année avec les autres thèmes.\\\bottomrule
\end{tabularx}
\begin{savoirfaire}[Une séance de 35 à 50 minutes]Relire une propriété et son exemple. Chercher deux exercices sans correction. Comparer les démarches et refaire une question ratée. Terminer par une phrase : « La prochaine fois, je dois penser à… ».\end{savoirfaire}
\begin{savoirfaire}[Revenir et mélanger]Reprendre une notion quelques jours plus tard, puis quelques semaines après. Mélanger un calcul, une lecture graphique et une preuve évite de dépendre seulement du titre d’un chapitre.\end{savoirfaire}
Pour un défi, prévoir 30 à 60 minutes, éventuellement en plusieurs séances.
Pour un bilan, respecter le temps choisi puis classer les erreurs par notion.
Le tableau de bord aide à décider ce qui doit être refait seul ou avec aide.

\clearpage
